\documentclass[11pt]{article}

\usepackage[margin=1in]{geometry}
\usepackage{lmodern}
\usepackage[T1]{fontenc}
\usepackage{microtype}
\usepackage{amsmath, amssymb, amsthm, mathtools, bm}
\usepackage{booktabs}
\usepackage{enumitem}
\usepackage{xcolor}
\usepackage{float}
\usepackage{algorithm}
\usepackage{algpseudocode}
\usepackage[colorlinks=true, linkcolor=blue!50!black, citecolor=blue!50!black,
            urlcolor=blue!50!black]{hyperref}

\newtheorem{proposition}{Proposition}
\theoremstyle{remark}
\newtheorem{remark}{Remark}

\newcommand{\x}{\bm{x}}
\newcommand{\y}{\bm{y}}
\newcommand{\uu}{\bm{u}}
\newcommand{\vv}{\bm{v}}
\newcommand{\eps}{\bm{\varepsilon}}
\newcommand{\zz}{\bm{z}}
\newcommand{\nn}{\bm{n}}
\newcommand{\cc}{\bm{c}}
\newcommand{\E}{\mathbb{E}}
\newcommand{\R}{\mathbb{R}}
\newcommand{\N}{\mathcal{N}}
\newcommand{\Id}{\bm{I}}
\newcommand{\Op}{\bm{E}}
\newcommand{\grad}{\nabla}
\newcommand{\PMI}{\operatorname{PMI}}
\newcommand{\ssb}{\sigma_{\mathrm{sb}}}
\newcommand{\code}[1]{\texttt{#1}}
\newcommand{\Dz}{D_{\varnothing}}
\newcommand{\Dg}{D_{\x_1}}
\newcommand{\Dy}{D_{\y}}
\newcommand{\Dgy}{D_{\x_1,\y}}

\title{Guided ImMAP ascent:\\[2pt]
       \large MAP with side information, one conditional expectation at a time}
\author{ImMAP project notes}
\date{November 2026}

\begin{document}
\maketitle

\begin{abstract}
We want the MAP estimate of a target image $\x_0$ given all available side information: a
guide image $\x_1$ (a paired scan, typically of another contrast), a measurement $\y$, or
both. Revealing the side information one variable at a time splits the log-posterior exactly
into the target prior plus one pointwise mutual information (PMI) term per variable, and each
PMI term is a likelihood. The split survives Gaussian smoothing, and by Tweedie's formula the
gradient of each term is the difference of two conditional expectations. So a single ascent
step handles any set of side information. The chain matters in practice because it says how
to \emph{build} the needed expectation from the ones we can estimate: e.g.\ $\E[\x_0\mid
\uu,\x_1,\y]$ from a guide-conditioned network plus a data step. I$^2$SB, I$^2$SB-Recon,
ImMAP2 and SDEdit are this one ascent with different conditioning sets, estimators and step
schedules.
\end{abstract}

\paragraph{Setting.} Real vectors in $\R^n$ (the complex case is identical with Wirtinger
gradients). Pairs $(\x_0,\x_1)$ have a joint law; the measurement is $\y = \Op\x_0 + \nn$,
$\nn\sim\N(\bm0,\sigma_n^2\Id)$, independent of $(\x_0,\x_1)$, so $\y$ and $\x_1$ are
conditionally independent given $\x_0$. $\uu = \x_0 + \sigma\eps$ with $\eps\sim\N(\bm0,\Id)$
independent of everything else; on the bridge it is the debiased state
$(\x_t - t\x_1)/(1-t)$ at $\sigma^2 = Vt/(1-t)$. For a set $\cc\subseteq\{\x_1,\y\}$ of
conditioning variables,
\[
  D_{\cc}(\uu) \coloneqq \E[\x_0\mid\uu,\cc]\quad\text{at the current noise level } \sigma .
\]

%=============================================================================================
\section{The objective}
%=============================================================================================

Given side information $\cc$, the goal is
\begin{equation}
  \hat\x_0 = \arg\max_{\x_0}\ \log p(\x_0\mid\cc),
  \qquad \cc\in\big\{\varnothing,\ \{\x_1\},\ \{\y\},\ \{\x_1,\y\}\big\}.
  \label{eq:map}
\end{equation}
We always condition on every independent source of information available. A variable that is
a function of others adds nothing: e.g.\ the zero-filled image $\Op^H\y$ is a function of
$\y$, so conditioning on it as well as $\y$ is conditioning on $\y$. A guide scan is not, so
$p(\x_0\mid\x_1,\y)\ne p(\x_0\mid\y)$ in general.

%=============================================================================================
\section{Decomposing along a chain}
%=============================================================================================

Reveal the variables of $\cc$ one at a time, in some order. With two variables, guide first:
\begin{align}
  \log p(\x_0\mid\x_1,\y)
  &= \log p(\x_0)
   + \underbrace{\log\frac{p(\x_0\mid\x_1)}{p(\x_0)}}_{\PMI(\x_0;\,\x_1)}
   + \underbrace{\log\frac{p(\x_0\mid\x_1,\y)}{p(\x_0\mid\x_1)}}_{\PMI(\x_0;\,\y\mid\x_1)}
  \label{eq:chain}\\
  &= \log p(\x_0) + \log p(\x_1\mid\x_0) + \log p(\y\mid\x_0) + \mathrm{const}.
  \label{eq:likelihoods}
\end{align}
\eqref{eq:chain} is an identity: the ratios telescope. Each PMI term measures how much the
newly revealed variable says about $\x_0$, given what was revealed before. In
\eqref{eq:likelihoods} each becomes a likelihood, by Bayes,
$\PMI(\x_0;\bm a\mid\bm b) = \log p(\bm a\mid\x_0,\bm b) - \log p(\bm a\mid\bm b)$, where the
last term does not depend on $\x_0$, and by the conditional independence
$p(\y\mid\x_0,\x_1) = p(\y\mid\x_0)$. So the objective is
\begin{center}
  \emph{target prior} $+$ \emph{guide likelihood} $+$ \emph{measurement likelihood},
\end{center}
whichever order the variables are revealed in, and dropping a variable drops its term. With
one variable the chain has one step, and with none \eqref{eq:map} is the prior alone.

The order does not change the objective. It changes the intermediate conditional
expectations that appear below, i.e.\ which estimators the gradient is assembled from:
guide-first uses $\Dz,\Dg,\Dgy$; measurement-first uses $\Dz,\Dy,\Dgy$.

%=============================================================================================
\section{Smoothing, and the gradient}
%=============================================================================================

As in ImMAP, we ascend a Gaussian smoothing of the objective while lowering the smoothing. Let
$p_\sigma(\uu\mid\cc) = \int\varphi_\sigma(\uu-\x)\,p(\x\mid\cc)\,d\x$, the law of the noisy
target given $\cc$, and
\[
  J_\sigma(\uu) \coloneqq \log p_\sigma(\uu\mid\cc),\qquad J_\sigma\to\log p(\cdot\mid\cc)
  \ \text{as}\ \sigma\to0 .
\]

\begin{proposition}[The chain survives smoothing; each term is a difference of expectations]
\label{prop:grad}
Along the chain $\varnothing\subset\cc_1\subset\dots\subset\cc_K = \cc$,
\begin{equation}
  J_\sigma(\uu) = \log p_\sigma(\uu)
  + \sum_{k=1}^K \log\frac{p_\sigma(\uu\mid\cc_k)}{p_\sigma(\uu\mid\cc_{k-1})},
  \qquad
  \sigma^2\grad J_\sigma(\uu) = \underbrace{\Dz + \sum_{k=1}^K\big(D_{\cc_k}-D_{\cc_{k-1}}\big)}
  _{=\ D_{\cc}(\uu)} -\ \uu ,
  \label{eq:grad}
\end{equation}
and each ratio is the PMI between the \emph{noisy} target $\uu$ and the $k$-th variable.
\end{proposition}
\begin{proof}
The telescoping is the same as in \eqref{eq:chain}, with $\uu$ in place of $\x_0$; Bayes applies
to $(\uu,\cc_k)$ because $\eps$ is independent of the side information. Tweedie's formula,
$\sigma^2\grad\log p_\sigma(\uu\mid\cc_k) = D_{\cc_k}(\uu) - \uu$, holds for any such $\cc_k$
(differentiate the convolution under the integral). Apply it to every term; the $-\uu$'s cancel
inside each ratio.
\end{proof}

\emph{What a variable contributes to the gradient is how far revealing it moves the
conditional expectation}, $D_{\cc_k} - D_{\cc_{k-1}}$; it vanishes when the variable is
uninformative given the others. Summed, the chain gives back $D_{\cc}$, so with an estimator of
$D_{\cc}$ itself nothing else is needed. The decomposition earns its keep when $D_{\cc}$ is
\emph{not} directly available: then one estimates a less-conditioned expectation and adds the
last increment separately. That is exactly what a data-consistency step does,
\[
  \Dgy = \Dg + \underbrace{(\Dgy - \Dg)}_{\text{measurement term}},
\]
with the increment estimated from the measurement model (Section~\ref{sec:estimators}).

The smoothed terms are not the unsmoothed likelihoods at a noisy image. The measurement term,
for instance, involves $p(\y\mid\uu,\x_1) = \int p(\y\mid\x)\,p(\x\mid\uu,\x_1)\,d\x$: the
likelihood averaged over the targets the noisy state still allows. It is weak at large $\sigma$
and sharpens to $p(\y\mid\x_0)$ as $\sigma\to0$.

%=============================================================================================
\section{The step}
%=============================================================================================

Gradient ascent on $J_\sigma$ with step size $h\sigma^2$ is, by \eqref{eq:grad},
\begin{equation}
  \boxed{\;\uu^+ = \uu + h\sigma^2\grad J_\sigma(\uu) = (1-h)\,\uu + h\,\vv,
  \qquad \vv = D_{\cc}(\uu).\;}
  \label{eq:step}
\end{equation}
By Tweedie's second-order identity, $-\grad^2 J_\sigma = (\sigma^2\Id -
\operatorname{Cov}[\x_0\mid\uu,\cc])/\sigma^4\preceq\sigma^{-2}\Id$, so every
$h\in(0,2)$ increases $J_\sigma$ (descent lemma), whatever $\cc$ is.

The step runs inside the ImMAP continuation~\cite{ks2021}: after each step, fresh noise is
injected so that the iterate is again a noisy target at the next, smaller level. With the
contraction $\sigma' = (1-\beta h)\sigma$, $\beta\in(0,1]$, an accurate step leaves noise
$(1-h)\sigma$, so the injection is
\begin{equation}
  \gamma = \sigma\sqrt{(1-\beta h)^2 - (1-h)^2}.
  \label{eq:inject}
\end{equation}
$\beta = 1$ injects nothing (a deterministic continuation toward the MAP estimate); $\beta<1$
samples the posterior.

\begin{algorithm}[H]
\caption{Guided ImMAP ascent for $\max_{\x_0}\log p(\x_0\mid\cc)$}
\label{alg:generic}
\begin{algorithmic}[1]
\Require side information $\cc$; an estimator of $D_{\cc}(\uu) = \E[\x_0\mid\uu,\cc]$ (directly,
         or assembled along a chain, \eqref{eq:grad}); start $\uu_0$ at level $\sigma_0$;
         steps $h_k\in(0,2)$; $\beta\in(0,1]$; $\sigma_{\min}$
\State $\uu\gets\uu_0$, $\sigma\gets\sigma_0$
\While{$\sigma>\sigma_{\min}$}
  \State $\vv \gets D_{\cc}(\uu)$
         \Comment{$=\uu + \sigma^2\grad J_\sigma(\uu)$}
  \State $\uu \gets (1-h)\,\uu + h\,\vv + \gamma\,\zz$, \ $\zz\sim\N(\bm0,\Id)$
         \Comment{step \eqref{eq:step}, injection \eqref{eq:inject}}
  \State $\sigma \gets (1-\beta h)\,\sigma$
\EndWhile
\State \Return $\vv$
\end{algorithmic}
\end{algorithm}

\noindent A natural start is the posterior mean, $D_{\cc}$ at $\sigma = \infty$, i.e.\
$\E[\x_0\mid\cc]$, re-noised to $\sigma_0$. An end-to-end reconstruction network (TGVN, LGGS)
that regresses $\E[\x_0\mid\y,\x_1]$ in one shot is exactly that start: the first step of the
ascent, without the rest.

%=============================================================================================
\section{Estimating the expectations}\label{sec:estimators}
%=============================================================================================

The algorithm consumes $D_{\cc}$ and is indifferent to its source:
\begin{itemize}[leftmargin=1.5em, itemsep=2pt]
  \item $\Dz$: a target-only denoiser, trained with the squared loss on $(\uu,\x_0)$.
  \item $\Dg$: the paired bridge regressor, trained with the squared loss on bridge states and
        called on $\x_t = (1-t)\uu + t\x_1$ at $t = \sigma^2/(\sigma^2+V)$.
  \item $\Dgy$, learned: a bridge regressor that also sees $(\y,\Op)$ and the guide (SB-TGVN,
        BridgeLGGS).
  \item $\Dgy$ (or $\Dy$) from $\Dg$ (or $\Dz$) by a \textbf{data step}: approximate the
        posterior given $\uu$ and the other variables by $\N(D, \varsigma^2\Id)$ around the
        estimator's output $D$; with the Gaussian likelihood, the posterior mean given $\y$ too
        is the proximal point
        \begin{equation}
          (\Op^H\Op + \lambda\Id)^{-1}(\Op^H\y + \lambda D),\qquad
          \lambda = \sigma_n^2/\varsigma^2,
          \label{eq:prox}
        \end{equation}
        solved by CG warm-started at $D$. Since the posterior variance is at most $\sigma^2$ on
        average ($\E\|\x_0 - D\|^2 = \mathrm{mmse}(\sigma)\le n\sigma^2$), the principled
        weight follows the noise level, $\lambda(\sigma) = \sigma_n^2/(c\sigma^2)$, $c\le1$.
\end{itemize}

%=============================================================================================
\section{Worked case: standard I\texorpdfstring{$^2$}{2}SB}\label{sec:i2sb}
%=============================================================================================

The source image $\x_1$ is the only side information; there is no measurement.

\paragraph{The objective.} $\cc = \{\x_1\}$, and the chain $\varnothing\subset\{\x_1\}$ has
one step:
\[
  \max_{\x_0}\ \log p(\x_0\mid\x_1)
  = \log p(\x_0) + \underbrace{\log p(\x_1\mid\x_0)}_{\PMI(\x_0;\,\x_1)\,+\,\mathrm{const}} .
\]
A target that is plausible on its own and consistent with this particular source. Smoothed
and differentiated (Proposition~\ref{prop:grad}),
\[
  \sigma^2\grad J_\sigma(\uu)
  = \underbrace{\big(\Dz - \uu\big)}_{\text{prior pull}}
  + \underbrace{\big(\Dg - \Dz\big)}_{\text{what } \x_1 \text{ adds}}
  = \Dg(\uu) - \uu .
\]
I$^2$SB never splits the two pulls: its network estimates their sum, $\Dg$, directly.

\paragraph{What the network estimates.} Training draws a pair $(\x_0,\x_1)$, a time
$t\sim\mathcal U[0,1]$ and a bridge state
\[
  \x_t = (1-t)\,\x_0 + t\,\x_1 + \ssb(t)\,\eps,\qquad \ssb(t)^2 = Vt(1-t),\quad V = 4\tau^2,
\]
and regresses $\x_0$ with the squared loss, so $D_\theta(\x_t,t,\x_1) = \E[\x_0\mid\x_t,\x_1]$.
Removing the known $\x_1$ part of the state turns it into a noisy target,
\[
  \uu = \frac{\x_t - t\,\x_1}{1-t} = \x_0 + \sigma\eps,\qquad \sigma^2 = \frac{Vt}{1-t},
\]
and $(\x_t,\x_1)\mapsto(\uu,\x_1)$ is invertible, so
$D_\theta(\x_t,t,\x_1) = \E[\x_0\mid\uu,\x_1] = \Dg(\uu)$ at $\sigma(t)$. \emph{Calling the
network at bridge time $t$ evaluates the gradient of the objective smoothed at
$\sigma(t)$}: $\sigma(t)^2\grad J_{\sigma(t)} = D_\theta - \uu$. The bridge time is the
smoothing level, from $\sigma = 0$ at $t=0$ to $\sigma=\infty$ at $t=1$.

\paragraph{One step, in both coordinates.} From time $t$ to $s<t$, with
$r = \sigma(s)/\sigma(t)$, i.e.\ $r^2 = \frac{s(1-t)}{t(1-s)}$, I$^2$SB takes the step
\eqref{eq:step} with $h = 1-r^2$ and the injection \eqref{eq:inject} with
$\beta = 1/(1+r)$:
\[
  \uu_s = \uu_t + (\sigma_t^2-\sigma_s^2)\,\grad J_{\sigma_t}(\uu_t) + \sigma_s\sqrt{1-r^2}\,\zz
        = r^2\,\uu_t + (1-r^2)\,\Dg(\uu_t) + \sigma_s\sqrt{1-r^2}\,\zz .
\]
The step size is the variance removed, and the injection tops the leftover noise
$r^2\sigma_t = r\sigma_s$ back up to exactly $\sigma_s$. Mapped back with
$\x_s = (1-s)\uu_s + s\,\x_1$, the $\x_1$ terms cancel and
\[
  \x_s = \tfrac{s}{t}\,\x_t + \big(1-\tfrac{s}{t}\big)\,\hat\x_0 + \sqrt{V\tfrac{s(t-s)}{t}}\;\zz,
  \qquad \hat\x_0 = D_\theta(\x_t,t,\x_1),
\]
which is I$^2$SB's reverse step: the bridge posterior between $\x_t$ and the predicted
$\hat\x_0$.

\paragraph{A run, from $\x_1$ to $\x_0$.} On a ladder $1 = t_0 > t_1 > \dots > t_K = 0$:
\begin{itemize}[leftmargin=1.5em, itemsep=2pt]
  \item \textbf{First step} ($t_0 = 1$, $\x_t = \x_1$, $\sigma = \infty$, $r = 0$, $h = 1$). The
        state carries no information beyond $\x_1$, so the network returns the posterior mean
        $\hat\x_0 = \E[\x_0\mid\x_1]$: a blurry average of every target consistent with $\x_1$.
        The step jumps to it and re-noises:
        $\x_{t_1} = (1-t_1)\,\E[\x_0\mid\x_1] + t_1\x_1 + \ssb(t_1)\zz$. This is a warm start
        at noise level $\sigma(t_1)$, not ``starting at $\x_1$''.
  \item \textbf{Middle steps} ($0<h<1$). Each moves the state the fraction $1-r^2$ toward the
        current estimate and re-injects noise. The fresh noise commits the run to one of the
        targets the estimate was averaging over; the next estimate, at lower $\sigma$, is
        sharper. Unrolled, the state is a precision-weighted average of all estimates so far.
  \item \textbf{Last step} ($s = 0$, $\sigma = 0$, $r = 0$, $h = 1$). No noise is injected and
        the output is the final estimate $\hat\x_0 = D_\theta(\x_{t_{K-1}},t_{K-1},\x_1)$.
\end{itemize}

\begin{table}[H]
\centering
\footnotesize
\begin{tabular}{@{}lcccccccc@{}}
\toprule
step $t\to s$ & $\sigma(t)$ & $\sigma(s)$ & $r$ & $h = 1-r^2$ & $\beta$ & injected (in $\uu$) &
  $\x_s$: weight on $\x_t$, $\hat\x_0$ & noise (in $\x$) \\
\midrule
$1\to0.75$    & $\infty$ & 0.658 & 0     & 1     & 1     & 0.658 & 0.75, 0.25 & 0.165 \\
$0.75\to0.5$  & 0.658    & 0.380 & 0.577 & 0.667 & 0.634 & 0.310 & 0.67, 0.33 & 0.155 \\
$0.5\to0.25$  & 0.380    & 0.219 & 0.577 & 0.667 & 0.634 & 0.179 & 0.50, 0.50 & 0.134 \\
$0.25\to0$    & 0.219    & 0     & 0     & 1     & 1     & 0     & 0, 1       & 0 \\
\bottomrule
\end{tabular}
\caption{Four uniform steps at $\tau = 0.19$ ($V = 0.144$). At the first step the weight
$0.75$ falls on $\x_t = \x_1$, i.e.\ $\x_{t_1} = 0.25\,\E[\x_0\mid\x_1] + 0.75\,\x_1 +
0.165\,\zz$. Uniform steps in $t$ give equal ratios $r$ in the middle of the ladder.}
\label{tab:i2sb_run}
\end{table}

\paragraph{A sample, not the MAP estimate.} $\beta = 1/(1+r)<1$ re-injects noise, so I$^2$SB
returns (approximately) a sample of $p(\x_0\mid\x_1)$. The approximation is the point estimate:
the exact reverse step would also carry the spread
$(1-r^2)^2\operatorname{Cov}[\x_0\mid\uu_t,\x_1]$ of the endpoint, which vanishes on fine
ladders. The same network gives the deterministic continuation toward the MAP estimate with
$\beta = 1$, i.e.\ $h = 1-r$ and no injection: the Euler step of the probability-flow ODE.
Keeping $h = 1-r^2$ and only dropping the noise (\code{deterministic=True}) is neither: the
state then lands at noise $r^2\sigma_t < \sigma_s$, below the level the next call is told.

%=============================================================================================
\section{The methods as instances}
%=============================================================================================

Every method below is Algorithm~\ref{alg:generic}. What differs is the conditioning set (what
is estimated), the estimator, and the hyperparameters: noise ladder, $h_k$, $\beta$, start.

\paragraph{I$^2$SB~\cite{liu2023}.} $\cc = \{\x_1\}$, $\vv = \Dg$, the paired bridge regressor;
fixed ladder $\sigma_k = \sigma(t_k)$ from $\infty$ to $0$ on NFE uniform bridge steps; with
$r_k = \sigma_{k+1}/\sigma_k$,
\begin{equation}
  h_k = 1 - r_k^2,\qquad \beta_k = \frac{1}{1+r_k},\qquad \gamma_k^2 = \sigma_{k+1}^2(1-r_k^2)
  \label{eq:i2sb}
\end{equation}
(Section~\ref{sec:i2sb}).

\paragraph{I$^2$SB-Recon and I$^2$SB-Inversion~\cite{wang2024}.} The guide is a fully sampled
scan of \emph{another contrast} (T1 for a T2 target; PD for PDFS here), and it is the bridge
start: $\x_1 = $ guide. The objective is the full posterior, $\cc = \{\x_1,\y\}$:
\[
  \log p(\x_0\mid\x_1,\y) = \log p(\x_0) + \log p(\x_1\mid\x_0) + \log p(\y\mid\x_0)
  + \mathrm{const}.
\]
The network only estimates $\Dg$ (it never sees $\y$); the measurement term is added by the data
step \eqref{eq:prox} from $\Dg$, with $\lambda = 7\cdot10^{-5}$ fixed and 7 CG iterations,
giving $\vv\approx\Dgy$. Since the guide is a separate scan, this term is genuinely informative:
the data step is the only place the measurement enters. Hyperparameters as I$^2$SB
\eqref{eq:i2sb}, with 80 steps. I$^2$SB-Inversion then walks the reconstruction back up the
bridge deterministically and runs the same ascent a second time from there.

\paragraph{ImMAP2 (\code{diffusion/immap.py}).} No guide: $\cc = \{\y\}$, $\vv\approx\Dy$, the
data step \eqref{eq:prox} applied to a target-only denoiser $\Dz$, with
$\lambda_k = \lambda_0\sigma_n^2(1+\sigma_k^2)/\sigma_k^2$ (posterior variance
$\sigma_k^2/(1+\sigma_k^2)$, the $\Pi$GDM rule) solved to tolerance. The ladder is self-paced,
$\sigma_k^2 = \tfrac1n\|\Dz(\uu_k)-\uu_k\|^2$ (slightly low on average:
$\E\|\Dz-\uu\|^2 = n\sigma^2 - \mathrm{mmse}$); $h_k = h_0k/(1+h_0(k-1))$ with $h_0 = 0.01$,
$\beta = 0.05$; start $\mathrm{mean}(\Op^H\y) + \zz$ at $\sigma = 1$. The implementation
returns the final state $\uu$, not $\vv$.

\paragraph{SDEdit.} $\cc = \varnothing$, $\vv = \Dz$: the source image enters only as the start
$\uu_0 = \x_1 + \sigma_0\zz$, not as conditioning.

\paragraph{A zero-filled bridge start.} If $\x_1 = \Op^H\y$ (our \code{sbunet} baseline), $\x_1$
is a function of $\y$, so $\cc = \{\x_1,\y\}$ is the same information as $\{\y\}$: the target is
the unguided posterior $p(\x_0\mid\y)$. The data step then adds no information in exact
arithmetic and only corrects the network's errors.

\begin{table}[H]
\centering
\footnotesize
\begin{tabular}{@{}p{3.6cm}p{2.0cm}p{4.3cm}p{4.6cm}@{}}
\toprule
method & conditions on & estimator of $D_{\cc}$ & ladder, $h$, $\beta$ \\
\midrule
SDEdit & nothing & $\Dz$ & free \\
ImMAP2 & $\y$ & data step on $\Dz$, $\lambda_k\propto(1+\sigma_k^2)/\sigma_k^2$ &
  self-paced; KS $h_k$; $\beta = 0.05$ \\
I$^2$SB & guide & $\Dg$ (bridge regressor) & fixed ladder; $h = 1-r^2$; $\beta = \frac1{1+r}$ \\
I$^2$SB-Recon / -Inversion & guide, $\y$ & data step on $\Dg$, $\lambda = 7\cdot10^{-5}$ &
  as I$^2$SB, 80 steps \\
SB-TGVN, BridgeLGGS (any start) & guide, $\y$ & learned $\Dgy$ (the net sees both) &
  as I$^2$SB \\
\code{sbunet} (zero-filled start) & $\y$ & $\Dg$ with $\x_1 = \Op^H\y$, + data step & as I$^2$SB \\
TGVN, LGGS end-to-end & guide, $\y$ & learned $\E[\x_0\mid\y,\x_1]$ & one step from
  $\sigma = \infty$ \\
\bottomrule
\end{tabular}
\caption{Every method is Algorithm~\ref{alg:generic}. The first two columns say what is
estimated; the last is how the ascent is run.}
\label{tab:methods}
\end{table}

Two consequences. First, I$^2$SB-Recon and ImMAP2 differ in what they condition on (the guide)
as well as in the estimator and hyperparameters, so comparing them conflates the three; the
table says which single change isolates each. Giving ImMAP2 the guide means replacing $\Dz$ by
$\Dg$; giving I$^2$SB-Recon ImMAP2's $\lambda_k$ isolates the data weight. Second, for SB-TGVN the
bridge start is a parameterisation, not an objective: whatever $\x_1$ is, the net conditions on
the guide and $\y$, so the three starts target the same posterior and differ only in the path.

\begin{thebibliography}{9}
\bibitem{ks2021}
Z.~Kadkhodaie, E.~P.~Simoncelli. Stochastic solutions for linear inverse problems using the
prior implicit in a denoiser. \emph{NeurIPS}, 2021.

\bibitem{liu2023}
G.-H.~Liu, A.~Vahdat, D.-A.~Huang, E.~A.~Theodorou, W.~Nie, A.~Anandkumar.
I$^2$SB: Image-to-Image Schr\"odinger Bridge. \emph{ICML}, 2023.

\bibitem{wang2024}
Y.~Wang et al. Guided MRI reconstruction via Schr\"odinger bridge. arXiv:2411.14269, 2024.
\end{thebibliography}

\end{document}
